% !TEX root = ../main.tex

\chapter{解耦式通用全身运动控制策略}

\section{上下肢解耦控制架构设计}

人形机器人在执行移动操作任务时，下肢负责维持稳定支撑和完成位姿调整，上肢负责执行抓取、挥拍、交互等精细动作。若将二者完全作为统一端到端控制问题求解，虽然理论上能够获得整体最优解，但在实际训练和部署中往往面临样本复杂度高、奖励设计困难、策略调试成本大以及迁移稳定性不足等问题。基于这一观察，本文提出上下肢解耦的通用全身控制架构，将移动稳定性与操作精度两个核心目标分别建模，并通过接口层实现统一协调。

在该架构中，下肢控制器接收速度命令、姿态目标和部分全身补偿信息，输出双腿关节动作以维持稳定 locomotion；上肢控制器接收来自遥操作模块或任务规划模块的末端目标、轨迹参考和姿态约束，并通过逆运动学和轨迹跟踪方法完成执行；躯干与状态管理模块则负责在上下肢之间同步关键姿态信息、监测质心偏移并触发必要的补偿策略。这样可以在保持系统模块化的同时，降低全身控制问题的求解难度。

从系统实现角度看，解耦架构还具有良好的可扩展性。对于纯遥操作场景，上肢接口可直接接收第 2 章重定向模块输出；对于自主任务场景，上肢接口也可无缝接收由视觉感知和高层规划器生成的末端目标。下肢控制器则可以在不同任务之间共享同一套基础移动能力，从而避免在每个任务中重复训练整套全身策略。


\section{上肢遥操作}

\subsection{遥操作数据接口}

上肢控制部分主要负责接收来自第 2 章重定向模块或任务规划模块的动作目标，并将其转化为可执行的上肢轨迹。为此，本文设计了统一的上肢数据接口，将输入抽象为末端位姿目标、关键中间路标点以及必要的姿态约束信息。该接口既支持人类示教产生的连续动作轨迹，也支持感知驱动任务中由规划器直接生成的目标位姿。

考虑到真实采集数据可能包含采样抖动、短时缺失或相邻关键帧间变化过快等问题，本文在接口层增加了低通滤波、插值与轨迹平滑处理机制。对于时序连续动作，可通过贝塞尔曲线或样条插值生成更平滑的参考轨迹；对于带有突然目标切换的任务，则通过限速和加速度约束避免对执行器造成过大冲击。

\subsection{逆运动学求解与跟踪}

在获得上肢末端目标后，系统需要进一步求解满足约束的关节角配置。本文采用解析方法与数值优化相结合的逆运动学策略：对结构规则、解析表达清晰的关节链可直接利用解析求解提高实时性；对于复杂姿态约束或接近奇异位形的情况，则使用数值迭代方法求得满足误差要求的解。通过两类方法的组合，可以兼顾求解速度与适用范围。

在跟踪执行阶段，系统不仅需要尽量减小末端误差，还需要保证关节运动连续、姿态变化平滑并避免超出机械极限。为此，本文在 IK 结果后增加关节限幅检查、速度限制和奇异点规避处理，并在必要时通过调整冗余自由度分配来改善解的可执行性。该设计使上肢控制器能够更稳定地服务于遥操作和自动执行两类场景。

\subsection{全身协调机制}

虽然上下肢在控制器层面被解耦，但在真实执行中二者并非完全独立。大范围上肢动作会导致机身质心偏移，快速下肢移动也可能影响上肢末端跟踪精度。因此，本文在系统层设计了全身协调机制，用于解决上下肢动作冲突和姿态补偿问题。

具体而言，当上肢执行大幅度伸展或快速挥动动作时，系统会根据当前机身姿态、支撑状态和期望动作方向，向下肢控制器发送附加姿态补偿需求，以减小重心偏移带来的不稳定因素；当下肢正在执行大幅转向或高速移动时，上肢控制器则适当降低高频动作幅度，优先保证末端轨迹的连续可达性。通过这种相互约束与信息共享机制，系统能够在保持模块化的同时实现全身协同。



\section{下肢强化学习 Locomotion 策略}

\subsection{策略训练环境配置}

下肢控制策略在 Isaac Gym 或 Isaac Lab 等大规模并行仿真环境中进行训练。仿真环境负责提供机器人动力学求解、接触模拟、地形变化以及外部扰动注入等功能，以支撑高频率的并行策略优化。训练过程中，机器人从不同初始姿态、速度指令和环境参数出发，不断学习如何维持平衡并实现期望移动行为。

为了提升策略在真实系统中的适应能力，本文在训练环境中引入域随机化机制，对摩擦系数、机身质量分布、执行器参数、传感噪声和地形高度扰动等因素进行随机采样。这样可以使策略在训练阶段接触更广泛的动力学条件，降低部署到真实机器人时因建模误差带来的性能退化风险。对于复杂地形或外部碰撞情形，也可通过额外扰动课程逐步增加训练难度。

\subsection{状态空间与奖励设计}

状态空间设计直接影响强化学习策略能否充分利用机器人本体信息并形成稳定步态。本文在下肢策略中使用的观测包括 IMU 姿态、机身角速度、各下肢关节角与关节速度、历史速度命令、足端接触信息以及必要的短时历史状态。通过这些输入，策略能够同时感知机器人当前姿态、运动趋势和外部支撑状态，从而做出合理动作决策。

动作空间采用下肢关节位置目标或位置增量的形式表示，覆盖左右腿共 12 个关键关节。相较于直接输出力矩控制指令，位置目标形式更便于与现有伺服执行接口兼容，也有助于提升训练稳定性。底层执行器再结合 PD 控制器或平台已有的关节控制模块完成实际跟踪。

奖励函数由多个子项组成，包括速度跟踪奖励、姿态稳定奖励、足端接触一致性奖励以及能量消耗惩罚等。其中，速度跟踪项鼓励策略准确跟随前进、后退、旋转和横向移动指令；平衡稳定项用于抑制过大的机身姿态偏差和跌倒；能量效率项则用于减少无效摆动和过大关节动作，以提升策略在真机上的可执行性。通过对不同奖励项进行加权组合，策略能够在稳定性与机动性之间取得平衡。

\subsection{网络结构与训练过程}

本文下肢策略采用基于 PPO 的强化学习算法进行训练。策略网络和价值网络均使用多层感知机结构，通过若干隐藏层对输入状态进行非线性特征提取，并输出各关节动作分布参数。PPO 算法在策略更新时具有较好的稳定性和实现成熟度，因此适合作为基础 locomotion 技能训练的核心优化方法。

在训练流程中，系统首先从较简单的速度命令和稳定站立任务开始，待策略具备基本平衡能力后，再逐步增加方向切换、复合速度命令和外部扰动。通过课程式训练，策略能够逐渐适应更复杂的控制目标。对于学习率、批次大小、折扣因子和裁剪系数等超参数，本文暂在附录中以占位方式保留，后续将结合最终实验配置补全。

\subsection{策略性能评估}

为了评价下肢 locomotion 策略的实用性，本文从速度覆盖范围、任务成功率和鲁棒性测试三方面进行分析。速度覆盖范围反映策略在不同方向上的可控运动区间；任务成功率用于统计在规定时间或路径约束下完成动作的比例；鲁棒性测试则关注策略在地形变化、外力扰动和指令切换条件下的稳定性。

结合当前提纲，本文重点关注前进、后退、原地旋转和左右平移四类基础运动模式。由于真机和仿真统计结果尚待整理，本节暂保留量化指标占位，如最大前进速度为 X\,m/s、旋转角速度为 X\,rad/s、任务成功率为 XX\% 等。后续补充正式实验结果后，可进一步将这些指标整理为统一表格，并与不同训练设置下的策略表现进行对比。

\section{Sim-to-Real 迁移与真机验证}

仿真训练获得的控制策略能否平稳部署到真实机器人，是评价系统工程价值的重要标准。为减小仿真与真实系统之间的差异，本文在控制部署阶段结合动力学参数辨识、执行器响应适配、传感器标定和控制频率对齐等手段，对策略输入输出接口进行统一修正。同时，通过在训练阶段引入域随机化，使策略提前适应一定范围内的模型偏差和环境变化。

在真机验证方面，本文计划围绕四类基础 locomotion 模式展开实验，包括前进、后退、原地旋转和左右平移，并考察不同速度指令下的稳定性、跟踪误差和能耗情况。对于上肢部分，则重点验证遥操作轨迹跟踪精度以及与下肢运动同时执行时的协调表现。对于尚未完成的统计结果，本文保留速度跟踪误差、关节力矩分布和单位任务能耗等指标占位，以待后续实验补齐。

总体而言，解耦式控制策略将遥操作、强化学习和任务接口三者连接起来，使本文的方法既能支撑基础移动技能，又能支持上层感知与规划模块的接入。这为第 4 章的任务级闭环验证奠定了基础。
